% !TEX root = template.tex

\begin{figure*}[!t]
    \centering
    \fcolorbox{gray}{white}{%
        \resizebox{\textwidth}{!}{
            \begin{tikzpicture}[x=1cm, y=1cm, >=Stealth]

    % ==========================================
    % MACROS PER DISEGNARE I BLOCCHI E I DETTAGLI
    % ==========================================

    % Macro per i tensori 3D Args: x, y, width, height, depth, color, label_text
    \newcommand{\drawbox}[7]{
        \pgfmathsetmacro{\dx}{#5*0.4}
        \pgfmathsetmacro{\dy}{#5*0.4}
        % Right side (depth)
        \draw[fill=#6!60!black, draw=black!60, line width=0.2mm] (#1+#3,#2) -- ++(\dx,\dy) -- ++(0,#4) -- ++(-\dx,-\dy) -- cycle;
        % Top side
        \draw[fill=#6!80!black, draw=black!60, line width=0.2mm] (#1,#2+#4) -- ++(\dx,\dy) -- ++(#3,0) -- ++(-\dx,-\dy) -- cycle;
        % Front face
        \draw[fill=#6, draw=black!80, line width=0.2mm] (#1,#2) rectangle ++(#3,#4);

        % 2. Centered Text Node inside the box
        \node[text width=#3 cm, align=center, font=\scriptsize, text=black, transform shape] at (#1+#3/2, #2+#4/2) {#7};
    }

    \newcommand{\slicebox}[7]{
        % #1=x, #2=y, #3=w, #4=h, #5=slice_d, #6=shift_d, #7=color
        \pgfmathsetmacro{\xshift}{#6*0.4}
        \pgfmathsetmacro{\yshift}{#6*0.4}
        \pgfmathsetmacro{\dx}{#5*0.4}
        \pgfmathsetmacro{\dy}{#5*0.4}

        \coordinate (FBL) at (#1+\xshift, #2+\yshift); % Front Bottom Left della fetta
        % Right side
        \draw[fill=#7!60!black, draw=black!80, line width=0.2mm] (#1+\xshift+#3, #2+\yshift) -- ++(\dx,\dy) -- ++(0,#4) -- ++(-\dx,-\dy) -- cycle;
        % Top side
        \draw[fill=#7!80!black, draw=black!80, line width=0.2mm]  (#1+\xshift, #2+\yshift+#4) -- ++(\dx,\dy) -- ++(#3,0) -- ++(-\dx,-\dy) -- cycle;
        % Front face
        \draw[fill=#7, draw=black!80, line width=0.2mm] (FBL) rectangle ++(#3,#4);
    }

    % 3. Macro per il tensore concatenato (Bianco + Rosso + Blu)
    \newcommand{\drawsplitbox}[6]{
        % #1=x, #2=y, #3=w, #4=h, #5=d, #6=label    
        % Retro: Ultimi 12 canali (Blu)
        \slicebox{#1}{#2}{#3}{#4}{0.2*#5}{0.9*#5}{myblue}
        % Centro: Successivi 256 canali (Rossi)
        \slicebox{#1}{#2}{#3}{#4}{0.4*#5}{0.45*#5}{myred}
        % Fronte: Primi 256 canali (Bianchi)
        \slicebox{#1}{#2}{#3}{#4}{0.4*#5}{0}{mywhite}
        % Testo sulla faccia frontale
        \node[text width=#3 cm, align=center, font=\scriptsize, text=black, transform shape] at (#1+#3/2, #2+#4/2) {#6};
    }

    % 4. Macro per il Residual Layer
    \newcommand{\resblock}[3]{
    % #1=startX, #2=startY, #3=endX
    \pgfmathsetmacro{\midX}{(#1+#3)/2}
    \pgfmathsetmacro{\plusY}{#2-0.4}
    \pgfmathsetmacro{\fY}{#2-0.8}

    % Box F(x) - Componenti verdi, testo nero
    \node[draw=mygreenLine, fill=mygreen!10, thick, text=black, minimum width=0.6cm, minimum height=0.4cm, font=\scriptsize, align=center] (fx) at (\midX, \fY) {$F(x)$};
    % Nodo Sommatore (+)
    \node[draw=mygreenLine, circle, inner sep=0pt, minimum size=0.25cm, thick, text=mygreenLine, font=\scriptsize] (plus) at (#3, \plusY) {$+$};

    % Rami con frecce più piccole e raffinate (scale=0.7)
    \draw[-{Stealth[scale=0.6]}, draw=mygreenLine, thick] (#1, #2) |- (fx.west);
    \draw[-{Stealth[scale=0.6]}, draw=mygreenLine, thick] (#1, #2) |- (plus.west);
    \draw[-{Stealth[scale=0.6]}, draw=mygreenLine, thick] (fx.east) -| (plus.south);
    \draw[-{Stealth[scale=0.6]}, draw=mygreenLine, thick] (plus.north) -- (#3, #2);
    }

    % Macro per il "piano roll" (i quadratini neri)
    \newcommand{\pianoroll}[2]{
        % Background box to contain the notes
        \draw[fill=white, draw=black, line width=0.2mm] (#1, #2) rectangle ++(1.5, 1.5);
        % Note blocks
        \fill[black] (#1+0.1, #2+1.0) rectangle ++(0.2, 0.1);
        \fill[black] (#1+0.3, #2+1.1) rectangle ++(0.2, 0.1);
        \fill[black] (#1+0.5, #2+0.8) rectangle ++(0.2, 0.1);
        \fill[black] (#1+0.7, #2+0.7) rectangle ++(0.2, 0.1);
        \fill[black] (#1+0.9, #2+0.9) rectangle ++(0.2, 0.1);
        \fill[black] (#1+1.1, #2+0.6) rectangle ++(0.2, 0.1);
        \fill[black] (#1+1.3, #2+0.5) rectangle ++(0.1, 0.1);
    }

    % Colori personalizzati per i blocchi
    \colorlet{myblue}{cyan!40}
    \colorlet{myred}{red!50}
    \colorlet{mygreen}{green!50}
    \colorlet{mywhite}{white}

    % Colori personalizzati (più scuri) per le linee e le frecce
    \colorlet{myblueLine}{cyan!70!black}
    \colorlet{myredLine}{red!70!black}
    \colorlet{mygreenLine}{green!70!black}

    % ==========================================
    % PARTE SINISTRA: CONDITIONER E GENERATOR
    % ==========================================

    % --- Titoli ---
    \node at (6, 9) {\Large \textbf{Generator with conditioner}};

    % ----------- MAIN PIPELINE -------------
    % Noise z (ruotato)
    \begin{scope}[rotate around={-6:(0,1)}]
        \drawbox{0}{5.7}{1.5}{0.3}{0}{mywhite}{\textbf{Noise $z$}};
    \end{scope}
    \draw[-{Stealth}, thick, gray] (1.5, 6.2) to[out=40, in=140] node[midway, above, text=black] {\scriptsize \textsf{project and reshape}} (3.0, 6.2);

    % First tensor: 524x1x2 (Utilizziamo la nuova \drawsplitbox)
    \drawsplitbox{3}{5.5}{0.5}{0.3}{1}{}
    \resblock{3.5}{5.3}{4.85}
    \draw[-{Stealth[scale=0.8]}, thick, gray] (4, 6.3) to[out=40, in=140] node[midway, above, text=black] {\scriptsize \textsf{transp conv}} (4.9, 6.3);

    % Second tensor: 524x1x4
    \drawsplitbox{4.5}{5.5}{1}{0.3}{1}{}
    \resblock{5.5}{5.3}{6.85}
    \draw[-{Stealth[scale=0.8]}, thick, gray] (6, 6.3) to[out=40, in=140] node[midway, above, text=black] {\scriptsize \textsf{transp conv}} (7, 6.3);

    % Third tensor: 524x1x8
    \drawsplitbox{6.5}{5.5}{1.5}{0.3}{1}{}
    \resblock{8}{5.3}{9.5}
    \draw[-{Stealth[scale=0.8]}, thick, gray] (8.5, 6.3) to[out=35, in=145] node[midway, above, text=black] {\scriptsize \textsf{transp conv}} (9.5, 6.3);

    % Fourth tensor: 524x1x16
    \drawsplitbox{9}{5.5}{2}{0.3}{1}{}
    \draw[-{Stealth[scale=0.8]}, thick, gray] (11, 6.3) to[out=60, in=160] node[midway, above, text=black] {\scriptsize \textsf{transp conv}} (11.7, 6.5);

    % Final Piano Roll
    \pianoroll{11.7}{5}
    \node at (12.45, 7) {\scriptsize Output};

    % ----------- CONDITION 1D -------------    
    \begin{scope}[rotate around={-6:(0,1)}]
        \drawbox{0}{8}{1}{0.3}{0}{myblue}{\textbf{Chord}};
    \end{scope}

    \draw[-{Stealth[scale=0.8]}, thick, myblueLine] (2, 8) -| (3.7, 6.3);
    \draw[-{Stealth[scale=0.8]}, thick, myblueLine] (2, 8) -| (5.5, 6.3);
    \draw[-{Stealth[scale=0.8]}, thick, myblueLine] (2, 8) -| (7.7, 6.3);
    \draw[-{Stealth[scale=0.8]}, thick, myblueLine] (2, 8) -| (10.2, 6.3);

    \node at (6, 8.2) {\scriptsize Extend};

    % ----------- CONDITION 2D -------------
    \pianoroll{0.5}{1.8}
    \draw[-{Stealth[scale=0.8]}, thick, gray] (2.1, 1.7) to[out=-35, in=-130] node[midway, below, text=black] {\scriptsize \textsf{conv}} (3, 2.2);

    \drawbox{3}{2.3}{2}{0.3}{0.4}{myred}{}
    \draw[-{Stealth[scale=0.8]}, thick, gray] (5, 2.2) to[out=-35, in=-145] node[midway, below, text=black] {\scriptsize \textsf{conv}} (6, 2.2);
    \drawbox{6}{2.3}{1.5}{0.3}{0.4}{myred}{}
    \draw[-{Stealth[scale=0.8]}, thick, gray] (7.5, 2.2) to[out=-35, in=-145] node[midway, below, text=black] {\scriptsize \textsf{conv}} (8.5, 2.2);
    \drawbox{8.5}{2.3}{1}{0.3}{0.4}{myred}{}
    \draw[-{Stealth[scale=0.8]}, thick, gray] (9.5, 2.2) to[out=-35, in=-145] node[midway, below, text=black] {\scriptsize \textsf{conv}} (10.5, 2.2);
    \drawbox{10.5}{2.3}{0.5}{0.3}{0.4}{myred}{}

    % |- arrows
    \draw[thick, myredLine] (10.75, 2.2) |- (11.40, 1.85);
    \draw[thick, myredLine] (9, 2.2) |- (11.55, 1.7);
    \draw[thick, myredLine] (6.75, 2.2) |- (11.7, 1.55);
    \draw[thick, myredLine] (4, 2.2) |- (11.85, 1.40);

    % |- arrows
    \draw[thick, myredLine] (11.40, 1.85) |- (3.25, 3.15);
    \draw[thick, myredLine] (11.55, 1.7) |- (5.15, 3.30) ;
    \draw[thick, myredLine] (11.7, 1.55) |- (7.15, 3.45);
    \draw[thick, myredLine] (11.85, 1.40) |- (10.15, 3.60);

    % final arrows
    \draw[-{Stealth[scale=0.8]}, thick, myredLine] (3.25, 3.15) -- (3.25, 5.3);
    \draw[-{Stealth[scale=0.8]}, thick, myredLine] (5.15, 3.30) -- (5.15, 5.3);
    \draw[-{Stealth[scale=0.8]}, thick, myredLine] (7.15, 3.45) -- (7.15, 5.3);
    \draw[-{Stealth[scale=0.8]}, thick, myredLine] (10.15, 3.60) -- (10.15, 5.3);

    % ==========================================
    % SEPARATORE CENTRALE
    % ==========================================
    \draw[line width=0.4mm, gray!40] (13.5, 0.5) -- (13.5, 8);

    % ==========================================
    % PARTE DESTRA: DISCRIMINATOR
    % ==========================================

    \node at (18.5, 9) {\Large \textbf{Critic}};

    % 1. Input Piano Roll
    \pianoroll{14.5}{5}

    % 2. Converging dashed lines for Conv Layers
    \draw[dashed, thick, black!70] (16.0, 6.5) -- (17.5, 5.8);
    \draw[dashed, thick, black!70] (16.0, 5.0) -- (17.5, 5.5);
    \node at (17.25, 6.7) {\scriptsize 4 conv steps};

    % 3. Conv Feature Map Block (256x1x1)
    \drawbox{17.5}{5.5}{0.3}{0.3}{0.6}{mywhite}{}

    % 4. Minibatch Standard Deviation (sigma)
    \draw[-{Stealth[scale=0.8]}, thick, gray] (17.65, 5.5) -- (17.65, 3.7) node[below, text=black] {\scriptsize $\sigma$};
    \node[anchor=south] at (16.15, 3.35) {\scriptsize std over minibatch};

    % 5. Chord Vector Condition
    \begin{scope}[shift={(14, 0)}] % Trasla il blocco in orizzontale, allineandolo al Critic
        \begin{scope}[rotate around={-6:(0,1)}] % Usa lo stesso identico asse di rotazione del lato sinistro
            \drawbox{0}{8}{1}{0.3}{0.1}{myblue}{\textbf{Chord}};
        \end{scope}
    \end{scope}

    % 6. Reshape and Concat Arrows
    % \draw[-{Stealth[scale=0.8]}, thick, gray] (18.2, 5.75) -- (20, 5.75) node[midway, below, text=black] {\scriptsize reshape \\ \& \\ concat};
    \draw[-{Stealth[scale=0.8]}, thick, gray] (18.2, 5.75) -- (20, 5.75) node[midway, below, text=black, align=center] {\scriptsize reshape \\ \scriptsize\& \\ \scriptsize concat};
    \draw[-{Stealth[scale=0.8]}, thick, myblueLine] (16, 8) -| (20.15, 7.1);
    \draw[-{Stealth[scale=0.8]}, thick, gray] (18, 3.5) -| (20.15, 3.9);

    % 7. Combined Flattened Vector (269 dims = 256 + 1 + 12)
    \drawbox{20}{4.0}{0.3}{3.0}{0}{mywhite}{}

    % 8. Fully Connected Convergence
    \draw[dashed, thick, black!70] (20.3, 7.0) -- (22.3, 5.7);
    \draw[dashed, thick, black!70] (20.3, 4.0) -- (22.3, 5.3);
    \node at (22.15, 6.5) {\scriptsize \textsf{2 linear steps}};

    % 9. Output Validity Node
    \drawbox{22.3}{5.3}{0.4}{0.4}{0}{mywhite}{V};
    \node at (22.5, 4.8) {\scriptsize Validity};

\end{tikzpicture}
        }
    }
    \caption{Model architecture: Generator with conditioner (left) and Critic (right).}
    \label{fig:model_architecture}
\end{figure*}

\section{Learning Framework}
\label{sec:learning_framework}

\subsection{Model Architecture}
\label{sec:model_architecture}

The noise given as inpute to the Generator $G$ is first processed throught two fully-connected layers, of $1024$ and $512$ neurons respectively, and then reshaped in a tensor of shape $(1, 256, 1, 2)$.
This is then followed by four convolutional layers, where at each stage a chord embedding $c \in \mathbb{R}^{12}$ is concatenated to the feature maps. The first three layers use filters of shape $(1, 2)$ and stride $2$, doubling the temporal dimesion at each step. To avoid vanishing gradients and ensure more stability, we integrated a Residual Block after each upsampling layer.
Finally, a transposed convolution with kernel $128 \times 1$ followed by a sigmoid activation procues the final bar of shape $128 \times 16$.

The Conditioner stricly follows the structure described in \cite{Yang-2017}. The aim is to enforce temporal continuity by conditioning the Generator on the previous bar. Its structure mirrors the one of the Generator. The output of each layer is concatenated with the corresponding Generator step along with its upsampling path.

Finally, the Critic is composed of four convolutional layers, with filters of shape $(4, 4)$ and stride $2$. The idea is to progressiveli extract features from the input bar, reducing spatial dimension from $128 \times 16$ to $1 \times 1$ with $256$ channels. Additnionally, we included a Minibatch Standard Deviation layer after the final convolutional stage. Its aim is to provide the network with a statistical information about the batch diversity, resulting in a scalar that is concatenated with the extracted features. These are then flattened, concatenated to the chord embedding and passed through a final fully connected layer, which outputs a single scalar value. Here we removed the Sigmoid activation, as it was not needed for the loss function we used.
\subsection{Wasserstein Loss and Gradient Penalty}
\label{sec:wasserstein_loss}

To address the vanishing gradient problem, we adopted the Wasserstein GAN with Gradient Penalty (WGAN-GP) as explained in \cite{Gulrajani-2017}.
The idea is to minimize the Wasserstein distance between the real distribution $P_r$ and the generated distribution $P_g$.
This distance is estimated via a 1-Lipschitz function $D$ (the Critic):
\begin{equation}
    W(P_r, P_g) = \sup_{\|D\|_L \leq 1} \mathbb{E}_{x \sim P_r}[D(x)] - \mathbb{E}_{\tilde{x} \sim P_g}[D(\tilde{x})]
\end{equation}
where $D(x)$ is the scalar output of the Critic, that can be seen as a score given to the generated music.

The losses for generator and Critic are computed as:
    
\begin{equation}
    L_D = \mathbb{E}_{\tilde{x} \sim P_g}[D(\tilde{x})] - \mathbb{E}_{x \sim P_r}[D(x)]
\end{equation}
\begin{equation}
    L_G = -\mathbb{E}_{\tilde{x} \sim P_g}[D(\tilde{x})]
\end{equation}

By minimizing $L_D$ the Critic learns to provide an increasingly acurate estimation of the Wasserstein distance. On the other hand, by maximizing $L_G$ the Generator learns to produce music that receives higher scores, effectively pushing $P_g$ closer to $P_r$.

Note that the Sigmoid activation is removed from the output of the Critic, allowing the scores to grow linearly and continuously providing useful gradients.


Regarding the Gradient Penalty, it is introduced to enforce the Lipschitz constraint on the Critic, acting as a regularization term on the gradient norm. The total loss of the Critic is then computed as:
\begin{equation}
    \begin{split}
        L = & \underbrace{\mathbb{E}_{\tilde{x} \sim P_g}[D(\tilde{x})] - \mathbb{E}_{x \sim P_r}[D(x)]}_{\text{Wasserstein Loss}}                                         \\
            & + \underbrace{\lambda \mathbb{E}_{\hat{x} \sim P_{\hat{x}}} \left[ \left( \|\nabla_{\hat{x}} D(\hat{x})\|_2 - 1 \right)^2 \right]}_{\text{Gradient Penalty}}
    \end{split}
\end{equation}

Note that $\hat{x}$ is a sampled taken from a random interpolation between real and fake samples. It can be seen as a random point along the line connecting the two distributions. This forces the Discriminator to keep small gradients along the trajectories connecting real and fake samples.

\subsection{Hyperparameter tuning}
The training was influenced mainly by three hyperparameters:
\begin{enumerate}
    \item \textbf{Gradient penalty coefficient ($\lambda$):} This parameter was never modified and was always set to $10.0$ as in the original paper.
    \item \textbf{Feature matching weight ($w_{fm}$):} As a validation metric we also used the Feature Matching Loss. The idea of this metric is to force the Generator to produce music that matches the statistics of the real data, as perceived by the Critic's intermediate layers. It is used to modify the total loss of the Generator:
    \begin{equation}
        \begin{split}
            L_G &= L_{adv} + w_{fm} \cdot L_{fm} \\
            \text{where } L_{adv} &= -\mathbb{E}_{\tilde{x} \sim P_g}[D(\tilde{x})] \\
            L_{fm} &= \| \mathbb{E}_{x \sim P_r} [f(x)] - \mathbb{E}_{\tilde{x} \sim P_g} [f(\tilde{x})] \|_2^2
        \end{split}
    \end{equation}
    where $f(\cdot)$ represents the feature extractions from the Critic's convolutional layers and $w_{fm}$ is the \textit{feature matching weight}.
    Varying this parameter allows us to control the balance between statistical accuracy and melodic variety.
    \item \textbf{Number of critic steps per generator step ($n_c$):} This parameter was fundamental after the change of loss function and critic structure, as it allowed us to balance the training of the two networks. We noticed that a value between 3 and 5 was optimal for the training, to allow the critic to learn enough to provide useful gradients to the generator, which otherwise tended to collapse
\end{enumerate}
\label{sec:hyperparameter_tuning}